% sec6_4.tex — Section 6.4: Estimating Autoregressions

Autoregressive processes (autoregressions) are popular in econometrics, not only
because they have a natural interpretation but also because they are easy to estimate.
This section considers the estimation of autoregressions. Estimation of
ARMA processes will be touched upon at the end.

\subsection*{Estimation of AR(1)}

Recall that an AR(1) is an $\text{MA}(\infty)$ with absolutely summable coefficients. So
if $\{\varepsilon_t\}$ is independent white noise (an i.i.d.\ sequence with zero mean and finite
variance), it is strictly stationary and ergodic (Proposition~6.1(d)). Letting $\mathbf{x}_t =
(1, y_{t-1})'$ and $\bm{\beta} = (c, \phi)'$, the AR(1) equation \eqref{eq:6.2.1} can be written as a regression
equation
\begin{equation}
  y_t = \mathbf{x}_t' \bm{\beta} + \varepsilon_t.
  \label{eq:6.4.1}
\end{equation}
We assume the sample is $(y_0, y_1, \ldots, y_n)$, $y_0$ inclusive, so that \eqref{eq:6.4.1} for $t = 1$
(which has $y_0$ on the right-hand side) can be included in the estimation. So the
sample period is from $t = 1$ to $n$.

We now show that all the conditions of Proposition~2.5 about the asymptotic
properties of the OLS estimator of $\bm{\beta}$ with conditional homoskedasticity are satisfied
here. First, obviously, linearity (Assumption~2.1) is satisfied. Second, as just
noted, $\{y_t, \mathbf{x}_t\}$ is jointly stationary and ergodic (Assumption~2.2). Third, since $y_{t-1}$
(the nonconstant regressor in $\mathbf{x}_t$) is a function of $\{\varepsilon_{t-1}, \varepsilon_{t-2}, \ldots\}$, it is independent
of the error term $\varepsilon_t$ by the i.i.d.\ assumption for $\{\varepsilon_t\}$. Thus
\begin{equation}
  \E(\varepsilon_t^2 \mid \mathbf{x}_t) = \sigma^2,
  \label{eq:6.4.2}
\end{equation}
that is, the error is conditionally homoskedastic (Assumption~2.7). Fourth, if
$\mathbf{g}_t \equiv \mathbf{x}_t \cdot \varepsilon_t = (\varepsilon_t, y_{t-1} \varepsilon_t)'$, $\{\mathbf{g}_t\}$ is a martingale difference sequence, as required in
Assumption~2.5, because

\medskip
\noindent first element of $\E(\mathbf{g}_t \mid \mathbf{g}_{t-1}, \mathbf{g}_{t-2}, \ldots)$
\begin{align*}
  &= \E(\varepsilon_t \mid \mathbf{g}_{t-1}, \mathbf{g}_{t-2}, \ldots) \\
  &= \E(\varepsilon_t \mid \varepsilon_{t-1}, \varepsilon_{t-2}, \ldots, y_{t-2}\varepsilon_{t-1}, y_{t-3}\varepsilon_{t-2}, \ldots) \\
  &= 0 \quad (\text{since $\{\varepsilon_t\}$ is i.i.d.\ and $y_{t-j}$ is a function of $(\varepsilon_{t-j}, \varepsilon_{t-j-1}, \ldots)$})
\end{align*}
and
\begin{align}
  &\text{second element of } \E(\mathbf{g}_t \mid \mathbf{g}_{t-1}, \mathbf{g}_{t-2}, \ldots) \notag \\
  &= \E(\varepsilon_t y_{t-1} \mid \mathbf{g}_{t-1}, \mathbf{g}_{t-2}, \ldots) \notag \\
  &= \E[\E(\varepsilon_t y_{t-1} \mid y_{t-1}, \mathbf{g}_{t-1}, \mathbf{g}_{t-2}, \ldots) \mid \mathbf{g}_{t-1}, \mathbf{g}_{t-2}, \ldots] \notag \\
  &\qquad \text{(by the Law of Iterated Expectations)} \notag \\
  &= \E[y_{t-1}\,\E(\varepsilon_t \mid y_{t-1}, \mathbf{g}_{t-1}, \mathbf{g}_{t-2}, \ldots) \mid \mathbf{g}_{t-1}, \mathbf{g}_{t-2}, \ldots] \notag \\
  &= 0 \quad (\text{since } \E(\varepsilon_t \mid y_{t-1}, \varepsilon_{t-1}, \varepsilon_{t-2}, \ldots, y_{t-2}\varepsilon_{t-1}, y_{t-3}\varepsilon_{t-2}, \ldots) = 0).
  \label{eq:6.4.3}
\end{align}
In particular, $\E(\varepsilon_t y_{t-1}) = 0$, so $\mathbf{x}_t$ is orthogonal to the error term (Assumption~2.3).
Finally, the rank condition about the moment matrix $\E(\mathbf{x}_t \mathbf{x}_t')$ (Assumption~2.4) is
satisfied because the determinant of
\begin{equation}
  \E(\mathbf{x}_t \mathbf{x}_t') =
  \begin{bmatrix} 1 & \mu \\ \mu & \gamma_0 + \mu^2 \end{bmatrix}
  \label{eq:6.4.4}
\end{equation}
is $\gamma_0 > 0$. This also means that $\E(\mathbf{g}_t \mathbf{g}_t')$ is nonsingular because under conditional
homoskedasticity $\E(\mathbf{g}_t \mathbf{g}_t') = \sigma^2 \E(\mathbf{x}_t \mathbf{x}_t')$. So the nonsingularity condition
in Assumption~2.5 is met.

Thus the AR(1) with an independent white noise process satisfies all the
assumptions of Proposition~2.5, so parts (a)--(c) of the propositions hold true here,
with $(c, \phi)'$ as $\bm{\beta}$, $(1, y_{t-1})'$ as $\mathbf{x}_t$, and $\sigma^2$ as the error variance. In particular, the
OLS estimate of the error variance given by
\[
  s^2 = \frac{1}{n-2} \sum_{t=1}^{n} e_t^2, \quad
  e_t \equiv y_t - \hat{c} - \hat{\phi}\,y_{t-1} \quad
  \text{where $(\hat{c}, \hat{\phi})$ are OLS estimates,}
\]
is consistent for $\sigma^2$.

\subsection*{Estimation of AR(\textit{p})}

Similarly for $\text{AR}(p)$, the autoregressive coefficients $(\phi_1, \ldots, \phi_p)$ can be consistently
estimated by regressing $y_t$ on the constant and lagged $y$'s, $(y_{t-1}, \ldots, y_{t-p})$.
We assume the sample is $(y_{-p+1}, y_{-p+2}, \ldots, y_0, y_1, \ldots, y_n)$, inclusive of $p$ lagged
values prior to $y_1$, so that the $\text{AR}(p)$ equation for $t = 1$ can be included and the
sample period is from $t = 1$ to $n$.

\begin{proposition}[Estimation of AR coefficients]
\label{prop:6.7}
Let $\{y_t\}$ be the $\mathrm{AR}(p)$ process
following \eqref{eq:6.2.6} with the stationarity condition \eqref{eq:6.1.18}. Suppose further that
$\{\varepsilon_t\}$ is independent white noise. Then the OLS estimator of $(c, \phi_1, \phi_2, \ldots, \phi_p)$
obtained by regressing $y_t$ on the constant and $p$ lagged values of $y$ for the sample
period of $t = 1, 2, \ldots, n$ is consistent and asymptotically normal. Letting
$\bm{\beta} = (c, \phi_1, \phi_2, \ldots, \phi_p)'$ and $\mathbf{x}_t = (1, y_{t-1}, \ldots, y_{t-p})'$, and $\hat{\bm{\beta}} =$ OLS estimate of
$\bm{\beta}$, we have
\[
  \avar(\hat{\bm{\beta}}) = \sigma^2\,\E(\mathbf{x}_t \mathbf{x}_t')^{-1},
\]
which is consistently estimated by
\[
  \widehat{\avar(\hat{\bm{\beta}})} = s^2 \cdot
  \left(\frac{1}{n} \sum_{t=1}^{n} \mathbf{x}_t \mathbf{x}_t'\right)^{-1},
\]
where $s^2$ is the OLS estimate of the error variance given by
\begin{equation}
  s^2 = \frac{1}{n - p - 1} \sum_{t=1}^{n}
  (y_t - \hat{c} - \hat{\phi}_1 y_{t-1} - \cdots - \hat{\phi}_p y_{t-p})^2.
  \label{eq:6.4.5}
\end{equation}
\end{proposition}

The proof is to verify the assumptions of Proposition~2.5, which can be done in the
same way as we did for AR(1). The only difficult part is to show that $\E(\mathbf{x}_t \mathbf{x}_t')$ is
nonsingular (Assumption~2.4), which is equivalent to requiring the $p \times p$ autocovariance
matrix $\Var(y_t, \ldots, y_{t-p+1})$ to be nonsingular (see Review Question~1(b)
below). It can be shown that the autocovariance matrix of a covariance-stationary
process is nonsingular for any $p$ if $\gamma_0 > 0$ and $\gamma_j \to 0$ as $j \to \infty$.\footnote{See, e.g., Proposition~5.1.1 of Brockwell and Davis (1991).}
So Assumption~2.4, too, is met.

We have not covered the maximum likelihood estimation of $\text{AR}(p)$, but for
those of you who are curious about the connection to ML, the OLS estimator is
numerically the same as the conditional Gaussian ML estimator and asymptotically
equivalent to the exact Gaussian ML estimator. See Section~8.7.

\subsection*{Choice of Lag Length}

All these nice results about estimation of autoregressions assume that the order
(the lag length) of autoregression, $p$, is known. How should we proceed if the
order $p$ is unknown? First of all, it is clear that Proposition~6.7, which is about
$p$-th order autoregressions, is applicable to autoregressions of lower order, say
$r < p$. That is, even if $\phi_r \neq 0$ but $\phi_{r+1} = \phi_{r+2} = \cdots = \phi_p = 0$, the proposition
remains applicable as long as $(\phi_1, \phi_2, \ldots, \phi_r)$ satisfies the stationarity condition.
In particular, the OLS estimates of $(\phi_{r+1}, \phi_{r+2}, \ldots, \phi_p)$ will converge to the true
value of zero. To put the same argument differently, suppose that the true order of
the autoregression is $p$ (so $\phi_p \neq 0$) and suppose that all we know about $p$ is that it
is less than or equal to some known integer $p_{max}$. By setting the $r$ in the argument
just given to $p$ and the $p$ to $p_{max}$, we see that Proposition~6.7 is applicable to
autoregressions with $p_{max}$ lags: if the $\{y_t\}$ is stationary $\text{AR}(p)$ and if we regress
$y_t$ on the constant and $p_{max}$ lags of $y$, then the OLS coefficient estimates will be
consistent and asymptotically normal.

It is natural to think that the true lag length may be estimated from this OLS
estimate of $(\phi_1, \ldots, \phi_{p_{max}})$. We consider two classes of rules for determining the
lag length. The first is the

\medskip
\noindent\textbf{General-to-specific sequential \textit{t} rule:} \; Start with an autoregression of $p_{max}$ lags.
If the last lag is significant at some prespecified significance level (say, 10 percent),
then end the procedure and set the lag length to $p_{max}$. Otherwise drop the
last lag, reestimate the autoregression with one fewer lag, and repeat the same
test. If this process continues until there is only one lag in the autoregression
and if that lag is insignificant, then set the lag length to 0 (no lags).

\medskip
This rule has the following interesting properties, to be illustrated in a moment. Since the
$t$-test is consistent, the lag length thus chosen will never be less than $p$ (the true lag
length) in large samples. However, the probability of overfitting (that the lag length
chosen is greater than the true lag length $p$) is not zero, even in large samples. That
is, if $\hat{p}$ is the lag length chosen by the sequential $t$ rule,
\begin{equation}
  \lim_{n \to \infty} \Pr(\hat{p} < p) = 0 \quad \text{and} \quad
  \lim_{n \to \infty} \Pr(\hat{p} > p) > 0.
  \label{eq:6.4.6}
\end{equation}

To illustrate, suppose $p = 2$ and $p_{max} = 3$. The sequential $t$ rule starts out
with three lags in the autoregression:
\[
  y_t = c + \phi_1 y_{t-1} + \phi_2 y_{t-2} + \phi_3 y_{t-3} + \varepsilon_t.
\]
We test the null hypothesis that $\phi_3 = 0$ at the 10 percent significance level. With
a probability of 10 percent, we reject the (true) null and set the lag length to~3, or
accept the null with a probability of 90 percent, in large samples. In the event of
accepting the null, we set $\phi_3 = 0$, estimate an autoregression with two lags, and test
the hypothesis that $\phi_2 = 0$. Since $\phi_2 \neq 0$ in truth (i.e., since $p = 2$), the absolute
value of the $t$-value on the OLS estimate of $\phi_2$ would be very large in large samples,
so we would never accept the (false) null that the second lag is zero. Therefore, in
this example, $\Pr(\hat{p} = 2) = 90\%$ and $\Pr(\hat{p} = 3) = 10\%$ in large samples.

There are two minor variations in the choice of the sample period with data
$(y_1, \ldots, y_n)$. One is to use the same sample period of $t = p_{max} + 1, p_{max} + 2, \ldots, n$
throughout. The other is to let the sample period be ``elastic'' by expanding it as
fewer and fewer lags are included. That is, when the autoregression being estimated
has $j$ lags, the sample period is $t = j + 1, j + 2, \ldots, n$.

The second rule for selecting the lag length uses either the \textbf{Akaike information
criterion} (\textbf{AIC}) or the \textbf{Schwartz information criterion} (\textbf{SIC}), also called the
\textbf{Bayesian information criterion} (\textbf{BIC}). This procedure sets the lag length to the $j$
that minimizes
\begin{equation}
  \log(SSR_j / n) + (j + 1) C(n) / n,
  \label{eq:6.4.7}
\end{equation}
over $j = 0, 1, \ldots, p_{max}$, where $SSR_j$ is the sum of squared residuals for the autoregression
with $j$ lags. In this formula, ``$j + 1$'' enters as the number of parameters
(the coefficients of the constant and $j$ lags). For the AIC, $C(n) = 2$ while for the
BIC, $C(n) = \log(n)$. A few comments about this information-criterion-based rule:
\begin{itemize}
\item As the number of parameters increases with $j$, the first term of the objective
function \eqref{eq:6.4.7} declines because the fit of the equation gets better, but the second
term increases. The information criterion strikes a balance between a better
fit and model parsimony.

\item Without an upper bound $p_{max}$, a ridiculously large value of $j$ might be chosen
(indeed, the information criterion achieves a minimum of $-\infty$ at $j + 1 = n$
because $SSR_j = 0$ for $j = n - 1$). In the present context where the true DGP is a
finite-order autoregression, a reasonable upper bound is the maximum possible
lag length.

\item Again there can be two ways to set the sample period. You can use the fixed
sample period of $t = p_{max} + 1, p_{max} + 2, \ldots, n$ throughout, in which case $SSR_j$
is the sum of $n - p_{max}$ squared residuals. Or you can use the ``elastic'' sample
period, in which case $SSR_j$ is the sum of $n - j$ squared residuals. In either case,
there is yet another variation, which is to replace the $n$ in the information criterion
\eqref{eq:6.4.7} by the actual sample size. So you replace the $n$ by $n - p_{max}$ if
the fixed sample is to be used. Based on simulations and some theoretical considerations,
Ng and Perron (2000) recommend the use of a fixed, rather than an
``elastic,'' sample period, with $n - p_{max}$ replacing $n$ in the information criterion.
\end{itemize}

Let $\hat{p}_{\text{AIC}}$ and $\hat{p}_{\text{BIC}}$ be the lag lengths chosen by the AIC and BIC, respectively.
Like the lag length selected by the general-to-specific $t$ rule, they are functions of
data, and hence are random variables. For them it is possible to show the following
(see, e.g., L\"utkepohl, 1993, Section~4.3):

\begin{enumerate}
\item[(1)] $\hat{p}_{\text{BIC}} \leq \hat{p}_{\text{AIC}}$ unless the sample size is very small. (If the fixed sample of
$t = p_{max} + 1, \ldots, n$ is used and ``$n - p_{max}$'' replaces $n$ in \eqref{eq:6.4.7}, this is true
for any finite sample with $n \geq p_{max} + 8$.) This is an algebraic result and holds
for any sample on $y$.

\item[(2)] Suppose that $\{y_t\}$ is stationary $\text{AR}(p)$ and $\{\varepsilon_t\}$ is independent white noise with
a finite fourth moment. Then
\begin{equation}
  \lim_{n \to \infty} \hat{p}_{\text{BIC}} = p, \quad
  \lim_{n \to \infty} \Pr(\hat{p}_{\text{AIC}} < p) = 0, \quad
  \lim_{n \to \infty} \Pr(\hat{p}_{\text{AIC}} > p) > 0.
  \label{eq:6.4.8}
\end{equation}
\end{enumerate}

That is, just like the general-to-specific sequential $t$ rule, the AIC rule has a positive
probability of overfitting. In contrast, the BIC rule is consistent. Furthermore,
Hannan and Deistler (1988, Theorem~5.4.1(c) and its Remark~1) have shown that
the consistency of $\hat{p}_{\text{BIC}}$ holds when the upper bound $p_{max}$ is made to increase at a
rate of $\log(n)$ (i.e., when it is set to $[c \log(n)]$ [the integer part of $c \log(n)$] for any
given $c > 0$). This means that you do not need to know an upper bound in the
consistent estimation by BIC of the order of an autoregression.

\subsection*{Estimation of VARs}

Estimation of VAR coefficients is equally easy. If $\mathbf{y}_t$ is $M$-dimensional, the $\text{VAR}(p)$
\eqref{eq:6.3.9} is an $M$-equation system and can be written as
\begin{equation}
  y_{tm} = \mathbf{x}_t' \bm{\delta}_m + \varepsilon_{tm} \quad (m = 1, 2, \ldots, M),
  \label{eq:6.4.9}
\end{equation}
where $y_{tm}$ is the $m$-th element of $\mathbf{y}_t$, $\varepsilon_{tm}$ is the $m$-th element of $\bm{\varepsilon}_t$, and
\[
  \underset{(Mp+1) \times 1}{\mathbf{x}_t} =
  \begin{bmatrix} 1 \\ \mathbf{y}_{t-1} \\ \mathbf{y}_{t-2} \\ \vdots \\ \mathbf{y}_{t-p} \end{bmatrix}, \quad
  \bm{\delta}_m =
  \begin{bmatrix} c_m \\ \bm{\phi}_{1m} \\ \bm{\phi}_{2m} \\ \vdots \\ \bm{\phi}_{pm} \end{bmatrix},
\]
$c_m = m$-th element of $\mathbf{c}$, $\bm{\phi}_{jm}' = m$-th row of $\bm{\Phi}_j$.
So the $M$ equations share the same set of regressors. It is easy to verify that $\mathbf{x}_t$
is orthogonal to the error term in each equation and that the error is conditionally
homoskedastic (the proof is very similar to the proof for $\text{AR}(p)$). So the
$M$-equation system with i.i.d.\ $\{\bm{\varepsilon}_t\}$ satisfies all the conditions for the multivariate
regression model of Section~4.5. Thus, efficient GMM estimation is very easy:
just do equation-by-equation OLS. The expressions for the asymptotic variance
and its consistent estimate can be obtained from Proposition~4.6 with $\mathbf{z}_{tm} = \mathbf{x}_t$ (so
$(4.5.13')$ on page~280 becomes $\frac{1}{n} \sum_t \mathbf{x}_t \mathbf{x}_t'$). If $\bm{\delta}$ is the $(Mp + 1)M$-dimensional
stacked vector created from $(\bm{\delta}_1, \ldots, \bm{\delta}_M)$ and if $\hat{\bm{\delta}}$ is the equation-by-equation OLS
estimate of $\bm{\delta}$, we have from (4.5.17)
\begin{equation}
  \widehat{\avar(\hat{\bm{\delta}})} = \hat{\bm{\Omega}} \otimes
  \left(\frac{1}{n} \sum_{t=1}^{n} \mathbf{x}_t \mathbf{x}_t'\right)^{-1},
  \label{eq:6.4.10}
\end{equation}
where
\begin{equation}
  \hat{\bm{\Omega}} \equiv \frac{1}{n - Mp - 1} \sum_{t=1}^{n} \hat{\bm{\varepsilon}}_t \hat{\bm{\varepsilon}}_t', \quad
  \hat{\bm{\varepsilon}}_t = \mathbf{y}_t - \hat{\mathbf{c}} - \hat{\bm{\Phi}}_1 \mathbf{y}_{t-1} - \cdots - \hat{\bm{\Phi}}_p \mathbf{y}_{t-p}.
  \label{eq:6.4.11}
\end{equation}

If we do not know the lag length $p$, we estimate it from data. In the information-criterion-based
rule, the lag length we choose is the $k$ that minimizes
\begin{equation}
  \log\left(\left|\frac{1}{n} \sum_{t=1}^{n} \hat{\bm{\varepsilon}}_t \hat{\bm{\varepsilon}}_t'\right|\right)
  + (k M^2 + M) \cdot \frac{C(n)}{n},
  \label{eq:6.4.12}
\end{equation}
over $k = 0, 1, \ldots, p_{max}$. Here, as in the univariate case, $C(n) = 2$ for the AIC and
$\log(n)$ for the BIC. Exactly the same results, listed as (1) and (2) above for the univariate
case, hold for the multivariate case (see, e.g., L\"utkepohl, 1993, Section~4.3).

\subsection*{Estimation of ARMA(\textit{p}, \textit{q})}

Letting
\begin{gather*}
  u_t \equiv \varepsilon_t + \theta_1 \varepsilon_{t-1} + \theta_2 \varepsilon_{t-2} + \cdots + \theta_q \varepsilon_{t-q}, \\
  \mathbf{z}_t = (1, y_{t-1}, y_{t-2}, \ldots, y_{t-p})', \quad
  \bm{\delta} = (c, \phi_1, \phi_2, \ldots, \phi_p)',
\end{gather*}
the univariate $\text{ARMA}(p, q)$ equation \eqref{eq:6.2.9} can be written as
\begin{equation}
  y_t = \mathbf{z}_t' \bm{\delta} + u_t.
  \label{eq:6.4.13}
\end{equation}
The moving-average component of ARMA creates two problems. First, obviously,
the error term $u_t$ is serially correlated (in fact, it is $\text{MA}(q)$). Second, since the
lagged dependent variables included in $\mathbf{z}_t$ are correlated with lags of $\varepsilon$ included in
the error term, the regressors $\mathbf{z}_t$ are not orthogonal to the error term $u_t$. That is,
serial correlation in the presence of lagged dependent variables results in biased
estimation. The second problem could be dealt with by using suitably lagged
dependent variables $(y_{t-q-1}, y_{t-q-2}, \ldots)$ as instruments. Regarding the first problem,
the method to be developed later in this chapter provides consistent estimates
of the coefficient vector in the presence of serial correlation. Given a consistent
estimate of $\bm{\delta}$, the MA parameters ($\theta$'s) can be estimated from the residuals. This
procedure, however, is not efficient.\footnote{The model yields infinitely many orthogonality conditions, $\E(u_t \cdot y_s) = 0$ for all $s \leq t - q - 1$. The GMM procedure can exploit only finitely many orthogonality conditions.}
We will not cover the topic of efficient estimation
of $\text{ARMA}(p, q)$ processes with Gaussian errors. Interested readers are
referred to, e.g., Fuller (1996, Chapter~8) or Hamilton (1994, Chapter~5).

\subsection*{Questions for Review}

\begin{enumerate}
\item (Avar of AR coefficients)
\begin{enumerate}[label=(\alph*)]
\item For AR(1), show that
\[
  \avar(\hat{\phi}) = \frac{\sigma^2}{\gamma_0} = 1 - \phi^2,
\]
where $\hat{\phi}$ is the OLS estimator of $\phi$. \textbf{Hint:} It is $\sigma^2$ times the $(2, 2)$ element
of the inverse of \eqref{eq:6.4.4}.

\item For $\text{AR}(p)$, let $\hat{\bm{\phi}}$ be the OLS estimate of $(\phi_1, \phi_2, \ldots, \phi_p)'$. Show that
\[
  \avar(\hat{\bm{\phi}}) = \sigma^2 \mathbf{V}^{-1},
\]
where $\mathbf{V}$ is the $p \times p$ autocovariance matrix
\[
  \mathbf{V} = \Var(y_t, \ldots, y_{t-p+1}) =
  \begin{bmatrix}
    \gamma_0 & \gamma_1 & \gamma_2 & \cdots & \gamma_{p-1} \\
    \gamma_1 & \gamma_0 & \gamma_1 & \cdots & \gamma_{p-2} \\
    \vdots & \ddots & \ddots & \ddots & \vdots \\
    \gamma_{p-2} & \cdots & \gamma_1 & \gamma_0 & \gamma_1 \\
    \gamma_{p-1} & \cdots & \gamma_2 & \gamma_1 & \gamma_0
  \end{bmatrix}.
\]
\textbf{Hint:} If $\mathbf{x}_t$ is as in Proposition~6.7, then
\[
  \E(\mathbf{x}_t \mathbf{x}_t') =
  \begin{bmatrix} 1 & \bm{\mu} \mathbf{1}' \\ \bm{\mu} \mathbf{1} & \mathbf{V} + \mu^2 \mathbf{1}\mathbf{1}' \end{bmatrix},
\]
where $\mathbf{1}$ is a $p \times 1$ vector of ones. Show that the inverse of this matrix is
\[
  \E(\mathbf{x}_t \mathbf{x}_t')^{-1} =
  \begin{bmatrix} * & * \\ * & \mathbf{V}^{-1} \end{bmatrix}.
\]
\end{enumerate}
\end{enumerate}
